\documentclass[12pt]{article}
\usepackage{amsmath, amssymb, geometry, booktabs}
\geometry{margin=1in}
\setlength{\parindent}{0pt}
\setlength{\parskip}{1em}

\title{Tutorial 4}
\author{ECON 101 -- Macroeconomics}
\date{Fall 2025}

\begin{document}
\maketitle

\section*{1. The Basic Multiplier in a Closed Economy (No Taxes, No Trade)}

\textbf{Consumption Function:}
\[
C = \bar{C} + cY
\]
where \( c \) is the marginal propensity to consume (MPC).

\textbf{Aggregate Expenditure:}
\[
AE = C + I + G = \bar{C} + cY + I + G
\]
Equilibrium occurs when \( Y = AE \).

\textbf{Multiplier (Closed Economy):}
\[
k = \frac{1}{1 - c}
\]

\textbf{Example Parameters:}
\[
\bar{C} = 200, \quad c = 0.75, \quad I = 150, \quad G = 250
\]

\textbf{Step 1: Solve for Equilibrium Income}
\[
Y = \bar{C} + cY + I + G
\]
\[
Y - 0.75Y = 200 + 150 + 250
\]
\[
0.25Y = 600 \quad \Rightarrow \quad Y = \boxed{2{,}400}
\]
\newpage
\textbf{Step 2: Apply a Government Spending Increase}
\[
\Delta G = +100, \quad k = \frac{1}{1 - 0.75} = 4
\]
\[
\Delta Y = k \cdot \Delta G = 4 \times 100 = \boxed{400}
\]
\[
Y' = 2{,}400 + 400 = \boxed{2{,}800}
\]

\textbf{Step 3: Round-by-Round Spending Illustration}
\[
100 + 0.75(100) + 0.75^2(100) + \dots = \frac{100}{1 - 0.75} = 400
\]
Each round induces additional consumption until total change in output equals \( \Delta Y = 400 \).



\section*{2. The Multiplier with Taxes and Imports (Open Economy)}

\textbf{Setup:}
\[
C = \bar{C} + c(1-t)Y, \quad NX = \bar{X} - mY
\]
\[
AE = \bar{C} + I + G + \bar{X} + [c(1-t) - m]Y
\]
Define \( A_0 = \bar{C} + I + G + \bar{X} \) (autonomous expenditure), and the slope of the AE function as \( z = c(1-t) - m \).

\textbf{Equilibrium:}
\[
Y = A_0 + zY \quad \Rightarrow \quad Y(1 - z) = A_0
\]
\[
Y = \frac{A_0}{1 - z}
\]

\textbf{Multiplier:}
\[
k = \frac{1}{1 - [c(1-t) - m]} = \frac{1}{1 - c(1-t) + m}
\]

\textbf{Example Parameters:}
\[
\bar{C} = 200, \; c = 0.75, \; I = 150, \; G = 250, \; \bar{X} = 20, \; t = 0.20, \; m = 0.15
\]

\textbf{Step 1: Compute the Slope \( z \)}
\[
z = c(1 - t) - m = 0.75(1 - 0.20) - 0.15 = 0.75(0.80) - 0.15 = 0.60 - 0.15 = \boxed{0.45}
\]

\textbf{Step 2: Compute Autonomous Spending \( A_0 \)}
\[
A_0 = \bar{C} + I + G + \bar{X} = 200 + 150 + 250 + 20 = \boxed{620}
\]

\textbf{Step 3: Find Equilibrium Income}
\[
Y = \frac{A_0}{1 - z} = \frac{620}{1 - 0.45} = \frac{620}{0.55} = \boxed{1{,}127.27}
\]

\textbf{Step 4: Compute the Multiplier}
\[
k = \frac{1}{1 - c(1-t) + m} = \frac{1}{1 - 0.75(0.80) + 0.15}
\]
\[
k = \frac{1}{1 - 0.60 + 0.15} = \frac{1}{0.55} = \boxed{1.8182}
\]

\textbf{Step 5: Fiscal Policy Shock}
\[
\Delta G = +100
\]
\[
\Delta Y = k \cdot \Delta G = 1.8182 \times 100 = \boxed{181.82}
\]
\[
Y' = 1{,}127.27 + 181.82 = \boxed{1{,}309.09}
\]

\textbf{Step 6: Spending Rounds Intuition (with Leakages)}
Each spending round induces a smaller response:
\[
100 + 0.45(100) + 0.45^2(100) + 0.45^3(100) + \dots = \frac{100}{1 - 0.45} = 181.82
\]

\textbf{Interpretation:}  
Taxes (\(t=0.20\)) and imports (\(m=0.15\)) create leakages that reduce the multiplier from \(4.0\) in a closed economy to approximately \(1.82\).  
This illustrates why fiscal policy has a smaller impact in open economies with taxation and trade.


\begin{center}
\rule{0.5\textwidth}{0.4pt}\\
\textit{End of Example}
\end{center}

\end{document}
